\section{Stable images and higher-defect conductor strata}
\label{sec:higher-defect}
The rank-two quotient in Theorem~\ref{thm:universal-hyperplane}
is specific to a hyperplane. In larger codimension the appropriate
object is a quotient algebra together with a subalgebra acting faithfully
on the quotient module. The rank of this action is an additional
stratification invariant. We keep its scheme structure throughout.

For a map $u$ of vector bundles write $I_j(u)$ for its ideal of
$j$-minors. Its \emph{exact-rank $s$ locus} means
$V(I_{s+1}(u))\setminus V(I_s(u))$, with this locally closed scheme
structure. On it the image and cokernel are vector bundles and commute
with arbitrary base change: after inverting a pivot, elementary
operations reduce the matrix to $\operatorname{diag}(I_s,0)$.
This is stronger than requiring constant rank on geometric points.

\begin{theorem}[Stable image and conductor stratification]
\label{thm:higher-defect-strata}
Let $B$ be a commutative finite locally free algebra of rank $d\ge2$
over a locally Noetherian $\C$-scheme $S$, and let $\mathcal A$ be the
universal bundle on $X=\Gr_S(k,B)^\circ$, with $1\le k\le d$.
Put $\mu_m:\Sym^m\mathcal A\to B_X$.
All Fitting ideals of $\operatorname{coker}\mu_m$ are independent of
$m\ge d-1$; Theorem~\ref{thm:codimension-stabilization} also gives
stability for $m\ge d-k+1$. We use $d-1$ below to retain a
uniform notation for the stable-image strata. On the exact-rank $s$ locus $X_s$ of $\mu_{d-1}$ there
is a canonical unital subalgebra bundle $E\subset B_{X_s}$ of rank $s$.
There is an invertible $E$-module $L\subset B_{X_s}$, also a subbundle
over $X_s$, such that
\begin{equation}\label{eq:stable-image-bundle}
 E=\mathcal O[a^{-1}\mathcal A],\qquad L=E\mathcal A=aE,
 \qquad \operatorname{im}\mu_m=L^m\quad(m\ge d-1)
\end{equation}
on any local unit cover $a\in\mathcal A\cap B^\times$.

Suppose $r=d-s>0$ and put $M=B/E$. The action
\[
 \theta:E\longrightarrow\mathcal E nd_{\mathcal O}(M)
\]
defines exact-rank strata $X_{s,t}$, with $1\le t\le r^2$.
On such a stratum the conductor and its quotient flag are
\begin{equation}\label{eq:conductor-flag}
 J=\ker\theta=\{b\in B:bB\subset E\},\qquad
 C=E/J\subset Q=B/J,
 \quad \rank C=t,\quad\rank Q=r+t.
\end{equation}
Here $J$ is an ideal subbundle of $B$, and $C\to\mathcal E nd(Q/C)$
is a subbundle injection. Each of these constructions commutes with
arbitrary base change within the specified stratum.

More generally, the functor of unital codimension-$r$ subalgebra
bundles with action of exact rank $t$ is equivalent to the functor of
flags $B\twoheadrightarrow Q\supset C$ of these ranks, where the
first map is a surjective morphism of unital commutative algebras,
$Q$ is locally free of rank $r+t$, $C\subset Q$ is a unital
subalgebra subbundle of rank $t$, $Q/C$ is locally free of rank $r$,
and the specified multiplication action
$C\to\mathcal E nd(Q/C)$ is a locally split subbundle injection. The inverse sends the flag to the
preimage of $C$.
\end{theorem}
\begin{proof}
Choose a unit $a$ smoothly locally, as in
Lemma~\ref{lem:generating-open}, and set $W=a^{-1}\mathcal A$.
Then $1\in W$. Lemma~\ref{lem:polarized-truncation} says that
$W^{d-1}=\mathcal O[W]$ and that $W^m=W^{d-1}$ for $m\ge d-1$,
over the coefficient ring itself. Multiplication by $a^m$ identifies
the original image with $a^m\mathcal O[W]$. The resulting cokernels
are locally isomorphic in these degrees, proving equality of every
Fitting ideal before any rank restriction.

On $X_s$, the pivot description makes $E=\mathcal O[W]$ a subbundle
with locally free quotient. If $b$ is a second unit in $\mathcal A$,
then $b/a\in E$ is a unit of $B$. Cayley--Hamilton for multiplication
by $b/a$, whose constant coefficient is a unit, puts its inverse
in $E$. Thus $\mathcal O[b^{-1}\mathcal A]\subset E$; reversing $a,b$
gives equality. This proves descent of $E$. The image $E\mathcal A$
is locally $aE$ and therefore descends to the stated invertible
$E$-module and subbundle $L$. It also proves
\eqref{eq:stable-image-bundle} universally.

Because $E$ is a subalgebra, it acts on $M=B/E$. Its action kernel
consists of the $e\in E$ with $eB\subset E$. An element $b$ of the
displayed conductor already lies in $E$, by multiplication by $1$;
hence the two kernels agree. If $jB\subset E$, then for $b,c\in B$
one has $(jb)c=j(bc)\in E$, so $J$ is an ideal of $B$.
On an exact-rank $t$ stratum, pivot reduction for $\theta$ makes
$J$ locally free of rank $s-t$ and $C$ a subbundle of
$\mathcal E nd(M)$. The exact sequence
$0\to C\to B/J\to M\to0$ proves that $Q$ is locally free of
rank $r+t$. All the kernels and images just used are split as
$\mathcal O$-modules locally; this proves their base-change assertions.
The unit acts as the identity on $M$, so $t\ge1$; the other bound
comes from $\rank\mathcal E nd(M)=r^2$.

Conversely, given a flag as in the statement, its inverse image $E$
is a subalgebra bundle and $B/E\simeq Q/C$. The kernel of its action
is precisely the kernel of $B\to Q$: a larger kernel would give
an element of $C$ annihilating $Q/C$, contrary to the subbundle
injection. Thus the conductor construction recovers the flag,
including over nonreduced test schemes. These mutually inverse
constructions respect all pullbacks within the exact-rank functors.
\end{proof}

The word ``stratum'' cannot be removed from the base-change assertion
for $J$. The multiplication cokernel is a universally defined
coherent sheaf on $X$; its conductor need not be a bundle on $X$.
Example~\ref{ex:conductor-rank-jump} below exhibits the distinction.

\begin{proposition}[The extreme-corank failure scheme]
\label{prop:extreme-corank}
On the Grassmannian of rank-$(d-r)$ subbundles $W\subset B$ containing
$1$, with $1\le r<d$, the unital subalgebra scheme is
\[
 V\!\left(\operatorname{Fitt}_{r-1}
       \operatorname{coker}(\Sym^m W\to B)\right),\qquad m\ge2.
\]
The defining ideal is independent of $m$, including its nonreduced
structure. Its conductor exact-rank $t$ subschemes are the quotient-flag
functors of Theorem~\ref{thm:higher-defect-strata}.
\end{proposition}
\begin{proof}
Since $1\in W$, the image $W^m$ contains $W$. Split $B=W\oplus M$
locally and remove the identity pivot columns. The residual
presentation has $r$ rows, so its $(r-1)$st Fitting ideal is the ideal
of all entries. In degree two these entries are precisely the
coefficients of $W\cdot W$ in $B/W$. They vanish if and only if $W$
is a subalgebra, on every test scheme. Modulo their ideal, every
higher product remains in $W$. Conversely each quadratic product is
an $m$-fold product by insertion of $1$'s. The entry ideals are equal.
The last assertion is the functorial equivalence just proved.
\end{proof}

For $r=1$ this is the usual zeroth-Fitting failure scheme. For $r>1$
it is the \emph{extreme-corank} locus, not the whole locus where
multiplication is nonsurjective. This distinction prevents a false
extension of the hyperplane theorem by changing only the quotient rank.

\subsection{The two quotient mechanisms in codimension two}

\begin{lemma}[Relative cyclic vectors]\label{lem:relative-cyclic-vector}
Let $M$ be a rank-two vector bundle, and let
$C\subset\mathcal E nd(M)$ be a unital rank-two algebra subbundle
over a $\C$-scheme. Then $M$ is an invertible $C$-module locally
on the base. More explicitly, write $C=\langle1,T\rangle$ and
$T=\left(\begin{smallmatrix}a&b\\c&d\end{smallmatrix}\right)$
after trivializing. The opens
\[
 D(c),\qquad D(b),\qquad D(c+d-a-b)
\]
cover the base and have cyclic vectors $(1,0)$, $(0,1)$, and
$(1,1)$, respectively. These opens and the resulting module
isomorphisms commute with base change.
\end{lemma}
\begin{proof}
On every geometric fibre, the subbundle injection makes $T$
nonscalar. If all three displayed functions vanish, then
$b=c=d-a=0$, contradicting that assertion. Their principal opens
therefore cover. The determinant of the map
$C\to M$, $q\mapsto qv$, is respectively $c$, $-b$, and
$c+d-a-b$ for the three chosen vectors. On the relevant open it
is a unit, so the map is a $C$-linear isomorphism. The same
invertible matrices prove the assertion after any pullback.
\end{proof}

\begin{theorem}[Codimension-two quotient flags]
\label{thm:codimension-two}
Let $E\subset B$ be a unital subalgebra bundle with $M=B/E$ of rank
two over a locally Noetherian $\C$-scheme. For its action $\theta$
one has $I_3(\theta)=0$ scheme-theoretically. The exact-rank-one
closed locus and the exact-rank-two open locus have the following
functorial descriptions.

On the first, $E$ is the inverse image of the scalar subalgebra in
a locally free rank-three quotient algebra of $B$.
On the second, $E$ is the inverse image of $C$ in a quotient
$B\twoheadrightarrow Q$, where $C$ is a rank-two algebra over the
base and $Q$ is a locally free rank-two algebra over $C$. In particular
$Q$ has rank four over the base. In both cases the kernel of
$B\to Q$ is the conductor, and these descriptions commute with
arbitrary base change on their respective strata.
\end{theorem}
\begin{proof}
Locally trivialize $M$. Since $2$ is invertible, write every matrix
as a scalar plus a trace-zero matrix. The trace-zero parts of the
commuting matrices $\theta(E)$ still commute. The bracket map
\[
 \bigwedge^2\mathfrak {sl}_2\longrightarrow\mathfrak {sl}_2
\]
is an isomorphism over $\C$: for the usual basis $H,U,V$ its values
are $2U,-2V,H$. Thus the wedge of any two trace-zero columns is zero
over the whole coefficient ring, not just over its residue fields.
Adjoining the scalar identity column proves $I_3(\theta)=0$.
The identity also gives $I_1(\theta)=\mathcal O$.

On the exact-rank-one locus the image is the scalar algebra. Theorem~
\ref{thm:higher-defect-strata} gives $C=\mathcal O$ and $\rank Q=3$.
Conversely the scalar line in any rank-three algebra has faithful
action on its rank-two quotient module: the action is ordinary scalar
multiplication. Its inverse image has exactly this conductor.

On the exact-rank-two locus, $C\subset\mathcal E nd(M)$ is a rank-two
subalgebra bundle. Locally write $C=\langle1,T\rangle$. Lemma~\ref{lem:relative-cyclic-vector}
gives an explicit cyclic-vector open cover and proves that $M$ is
locally free of rank one as a $C$-module. The exact sequence of $C$-modules
$0\to C\to Q\to M\to0$ consequently splits locally, so $Q$ is
locally free of rank two over $C$.

Conversely, if $Q$ is a locally free rank-two $C$-algebra, its unit
is unimodular as a $C$-section. Hence $Q/C$ is locally free of rank
one over $C$. Its regular $C$-action is faithful and is a subbundle
injection into its endomorphism bundle over the original base
(check on geometric fibres and use a pivot). The quotient-flag
equivalence therefore gives exactly the second stratum. All arguments
are functorial and use split modules on the indicated strata.
\end{proof}

\begin{remark}\label{ex:conductor-rank-jump}
Consider $R=\C[s]$, $B=R[z]/(z^5)$, and the direct summand
\[
 E=\langle1,sz^2+z^3,z^4\rangle_R.
\]
It is a subalgebra since $(sz^2+z^3)^2=s^2z^4$. Modulo $E$ take the
basis $\bar z,\overline{z^2}$; then $\overline{z^3}=-s\overline{z^2}$.
The action of $sz^2+z^3$ sends $\bar z$ to $-s^2\overline{z^2}$ and
$\overline{z^2}$ to zero, while $z^4$ acts as zero. Thus the
exact-rank-one subscheme is \emph{$s^2=0$}, not merely $s=0$.
Over $R$ the conductor is $Rz^4$; at the reduced fibre $s=0$ it is
$\langle z^3,z^4\rangle$. Quotient rank drops from four to three.
This simultaneously records the nonreduced rank stratum and the failure
of conductor formation to commute with unrestricted specialization.
The fibre subalgebras belong to the classical codimension-two
polynomial-algebra classification of \cite[Theorem 39]{GLTU}; the
calculation here retains the relative action and its determinantal
thickening.
\end{remark}
